\documentclass[11pt]{article}
\usepackage[T1]{fontenc}
\usepackage[utf8]{inputenc}
\usepackage[spanish]{babel}
\usepackage{amsmath,amssymb,amsthm}
\usepackage{booktabs}
\usepackage[margin=2.5cm]{geometry}

\newtheorem*{resultado}{Resultado}

\title{La prueba de robustez del reparto de la pendiente de accidentes}
\author{Notas de Claude para Rodrigo}
\date{17 de noviembre de 2026}

\begin{document}
\maketitle

\begin{abstract}
Explico qué era la prueba de robustez $\alpha_1=\alpha_1^G/2$, por qué se propuso y qué
descubrí al revisarla con cuidado: con nuestra especificación, cambiar el reparto de la
pendiente equivale a cambiar un poco $u_1$, la pendiente de la utilidad en la edad. Por eso
esa prueba casi no prueba nada. Corrijo lo que te dije en el chat (que ``reparar vale más''
en esa variante): no es así. Al final está la pregunta que necesito que contestes.
\end{abstract}

\section{Dos fuentes de deterioro}

En la tesis la probabilidad de accidente es
\[
s(j,a,w)=\Lambda\big(\alpha_{0j}+\alpha_{1j}\,a+w\big),
\qquad
w_{t+1}=w_t+\delta_j-\kappa\,r_t+\eta_t .
\]
El log-momio del accidente tiene dos partes que crecen con los años:
\begin{itemize}
  \item \textbf{Envejecimiento} $\alpha_{1j}a$: sube $\alpha_{1j}$ cada año y la reparación no
  lo toca.
  \item \textbf{Desgaste} $w$: sube $\delta_j$ cada año en que no se repara; una reparación lo
  baja $\kappa$.
\end{itemize}

Gillingham et al.\ solo tienen la curva $\Lambda(\alpha^G_{0j}+\alpha^G_{1j}a)$: una sola
pendiente $\alpha^G_{1j}$ (0.1804 en light brown, 0.2020 en heavy brown). Esa pendiente se
estimó con coches reales, que se reparan, así que ya mezcla las dos partes.

\section{El problema: una ecuación y dos incógnitas}

Si un coche se repara con una tasa $\bar p$, su log-momio sube en promedio
\[
\alpha_{1j}+\delta_j-\kappa\bar p \quad\text{por año.}
\]
Para que nuestra economía se parezca a la de Gillingham pedimos que esto sea igual a
$\alpha^G_{1j}$:
\begin{equation}\label{eq:cal}
\alpha_{1j}+\delta_j-\kappa\bar p=\alpha^G_{1j}.
\end{equation}
Es una sola ecuación para dos números, $\alpha_{1j}$ y $\delta_j$. Hay muchas formas de repartir
la pendiente. Las escribo con un parámetro $\lambda\in[0,1]$, la fracción de la pendiente de
Gillingham que se le atribuye al desgaste reparable:
\[
\alpha_{1j}=(1-\lambda)\,\alpha^G_{1j},
\qquad
\delta_j=\kappa\bar p+\lambda\,\alpha^G_{1j}.
\]

\begin{center}
\begin{tabular}{lccc}
\toprule
 & $\lambda$ & $\alpha_1$ (light brown) & $\delta$ (light brown)\\
\midrule
Caso base (el que está en el código y la tesis) & 0 & 0.1804 & 0.075\\
Robustez propuesta en octubre & 1/2 & 0.0902 & 0.165\\
\bottomrule
\end{tabular}
\end{center}

La idea de la robustez era: como el reparto no lo dicen los datos de Gillingham, sino que lo
elegimos nosotros, hay que revisar que los resultados del Monte Carlo no dependan de esa
elección.

\section{Lo que encontré: el reparto casi no cambia la economía}

Primero un ejemplo. Tomo un coche light brown y calculo cuánto sube su log-momio por año
(sin el choque $\eta$):

\begin{center}
\begin{tabular}{lcc}
\toprule
 & $\lambda=0$ & $\lambda=1/2$\\
\midrule
Nunca se repara: $\alpha_1+\delta$ & $0.1804+0.075=0.255$ & $0.0902+0.165=0.255$\\
Se repara cada año: $\alpha_1+\delta-\kappa$ & $0.255-0.25=0.005$ & $0.255-0.25=0.005$\\
\bottomrule
\end{tabular}
\end{center}

Es lo mismo. La razón es que $\alpha_1+\delta$ no depende de $\lambda$, y lo que hace la
reparación ($-\kappa$) tampoco. Lo único que cambia es \emph{qué número llamamos} $w$.

\begin{resultado}
Tomemos la misma historia de reparaciones y de choques $\eta$. Si $w^0$ es el desgaste en el
caso base y $w^\lambda$ el de la variante, entonces
\[
w^\lambda_a=w^0_a+\lambda\,\alpha^G_{1j}\,a .
\]
Por lo tanto:
\begin{enumerate}
  \item La probabilidad de accidente es idéntica:
  $\alpha_0+(1-\lambda)\alpha^G_1a+w^\lambda_a=\alpha_0+\alpha^G_1a+w^0_a$.
  \item El efecto de reparar y el precio de reparar son idénticos.
  \item La única diferencia está en la utilidad, porque $u_w\,w$ entra en ella:
  \[
  u_1a+u_w\,w^\lambda_a=\big(u_1+u_w\lambda\alpha^G_{1j}\big)a+u_w\,w^0_a .
  \]
\end{enumerate}
Es decir, el modelo con reparto $\lambda$ es el modelo base con
$u_1\to u_1+u_w\lambda\alpha^G_{1j}$, después de renombrar $w$. Si $u_w=0$, son exactamente
el mismo modelo, salvo la discretización del grid.
\end{resultado}

\noindent
\textit{Demostración.} Los coches nuevos nacen en $w=0$ en los dos casos. Cada año
$w^\lambda$ sube $\delta^\lambda=\delta^0+\lambda\alpha^G_1$ y $w^0$ sube $\delta^0$, con la
misma reparación y el mismo $\eta$. Así que la diferencia crece exactamente
$\lambda\alpha^G_1$ por año. Los puntos 1 a 3 se siguen de sustituir. \qed

\medskip
Con los números de la calibración ($u_w=-0.2$, $\lambda=1/2$), el cambio en $u_1$ es
$-0.2\times0.09\approx-0.018$ en light brown. La pareja pobre pasa de $u_1=-0.146$ a $-0.164$:
la utilidad cae un poco más rápido con la edad. Las decisiones de reparar, los precios por
historia del coche y la información de la mezcla para identificar $r$ son los mismos.

\paragraph{Corrección de lo que te dije.} En el chat escribí que en la variante ``una parte
mayor del deterioro se puede reparar y reparar vale más''. Es falso: lo que una reparación
quita ($\kappa$) es igual en los dos casos, y lo que se acumula sin reparar ($\alpha_1+\delta$)
también. La variante tampoco cambia lo que vale reparar. Solo mueve $u_1$ un poco y ensancha
el rango que necesita el grid de $w$ (por eso se hablaba de $W\approx100$).

\paragraph{Algo útil para la tesis.} Esto también aclara la nota al asesor de la
calibración (l.~757). El reparto entre $\alpha_1$ y $\delta$ es casi una normalización. Fija
dónde ponemos el cero de $w$ en cada edad y solo importa a través de $u_w$. Se puede decir en
una frase y deja de ser una elección incómoda.

En la estimación $\alpha_1$ y $\delta$ sí se separan, porque $w$ se observa: $\alpha_1$ sale de
cómo cambian los accidentes con la edad para un mismo $w$, y $\delta$ de la transición de $w$.
Lo que no importa es cuál reparto usamos para \emph{generar} los datos.

\section{La pregunta}

Como esa robustez equivale a mover $u_1$ un 12\%, correrla cuesta un Monte Carlo completo
(y con un grid más grande) y casi no informa. Las robusteces que sí ponen a prueba la
identificación de la reparación no observada son las que mueven la información de la mezcla
o de la variable excluida:

\begin{enumerate}
  \item[(a)] \textbf{Mezcla menos separada:} $\kappa$ más chico respecto a $\sigma_\eta$ (por
  ejemplo, $\kappa=0.15$ con $\sigma_\eta=0.15$, en vez de $0.25$). Las densidades de
  ``reparó'' y ``no reparó'' se enciman más, y es más difícil inferir $r$.
  \item[(b)] \textbf{Menos variación en el precio de reparar:} regímenes
  $\zeta\in\{-0.1,0,0.1\}$ en vez de $\pm0.3$. La variable excluida mueve menos la
  probabilidad de reparar.
  \item[(c)] \textbf{Sin utilidad del desgaste:} $u_w=0$. La reparación solo sirve para alargar
  la vida y subir el precio de reventa, no para la comodidad.
\end{enumerate}

\noindent\textbf{Pregunta:} ¿qué hacemos con la sección de robustez?
\begin{enumerate}
  \item[\textbf{A.}] Quitar la robustez del reparto $\lambda$ y poner en la tesis una nota al pie
  con el resultado de la sección 3. No implementar nada más por ahora.
  \item[\textbf{B.}] Lo mismo que A, y además implementar algunas de (a), (b) o (c) como
  robusteces. Cada una es un Monte Carlo completo en GPU. Dime cuáles.
  \item[\textbf{C.}] Correr de todos modos la de $\lambda=1/2$, como confirmación numérica del
  resultado.
\end{enumerate}

\noindent Mi recomendación es \textbf{B con (a) y (b)}. Son las que contestan la pregunta
central de la tesis: cuánta información hace falta para recuperar una decisión que no se ve.
La (c) es más bien sobre el canal económico de la reparación.

\end{document}
